\documentclass[12pt, aspectratio=169]{beamer}
\usetheme{Madrid}
\usecolortheme{default}

\usepackage{cmap}					    %%% PDF search
\usepackage{mathtext} 				    %%% cyrillic in math mode

\usepackage[T2A]{fontenc}
\usepackage[utf8]{inputenc}
\usepackage[russian]{babel}
\usepackage{amsmath, amssymb, amsfonts}
\usepackage{booktabs}
\usepackage{array}
\usepackage{multicol}

\usepackage{pgfplots}
\pgfplotsset{compat=1.18}
\usepackage{breqn}

\usepackage{hyperref}
\usepackage{graphicx}
\usepackage{pgfplots}
\usepackage{pgfplotstable}
\usepgfplotslibrary{groupplots}
\usetikzlibrary{calc}
\pgfplotsset{compat=1.18}

\usepgfplotslibrary{external}
\pgfplotsset{table/search path={., ./attachments}}
\tikzexternalize[prefix=tikz/]

\graphicspath{{attachments/}}

% Temporary workaround
% See https://github.com/matplotlib/matplotlib/issues/27907
\def\mathdefault#1{#1}
\everymath=\expandafter{\the\everymath\displaystyle}
\IfFileExists{scrextend.sty}{
  \usepackage[fontsize=10.000000pt]{scrextend}
}{
  \renewcommand{\normalsize}{\fontsize{10.000000}{12.000000}\selectfont}
  \normalsize
}


\setbeamertemplate{navigation symbols}{}

% ============================================================
\makeatletter
\setbeamertemplate{footline}
{
  \leavevmode%
  \hbox{%
  \begin{beamercolorbox}[wd=.25\paperwidth,ht=2.25ex,dp=1ex,left]{author in head/foot}%
    \usebeamerfont{author in head/foot}\quad\insertauthor
  \end{beamercolorbox}%
  \begin{beamercolorbox}[wd=.5\paperwidth,ht=2.25ex,dp=1ex,center]{title in head/foot}%
    \usebeamerfont{title in head/foot}\inserttitle
  \end{beamercolorbox}%
  \begin{beamercolorbox}[wd=.25\paperwidth,ht=2.25ex,dp=1ex,right]{date in head/foot}%
    \usebeamerfont{date in head/foot}\insertshortdate{}\hspace*{2em} \insertframenumber{} / \inserttotalframenumber\hspace*{2ex}%
  \end{beamercolorbox}}%
  \vskip0pt%
}
\makeatother

%=====================================================
% Title page definition
% ============================================================
\title{Модифицированный метод Лакса-Фридрихса для уравнения переноса}
\author{Уколов Иван.}
\institute{
  МФТИ\\
  Кафедра Математического моделирования и прикладной математики.\\
  Научный руководитель: И.В. Попов.
}
\date{г. Москва, 2026 г.}



%=====================================================
\begin{document}


%=====================================================
% Title page
%=====================================================
\begin{frame}
\titlepage
\end{frame}


%=====================================================
\begin{frame}{Постановка задачи}
Рассматривается одномерное линейное уравнение переноса
\[
\frac{\partial u}{\partial t} + a\,\frac{\partial u}{\partial x} = 0,\qquad a = \mathrm{const} > 0,
\]
с начальным условием
\[
u(x,0)=f(x),\quad x\in[0,L].
\]
Точное решение — перенос профиля без искажений:
\[
u(x,t)=f(x-at).
\]

\textbf{Цель:} построить явную разностную схему повышенного (третьего) порядка, сохраняющую монотонность на разрывах.
\end{frame}


%=====================================================
\begin{frame}{Модификация схемы Лакса--Фридрихса}
Классическая схема Лакса--Фридрихса:
\[
u_i^{n+1} = \frac{1}{2}(u_{i+1}^n + u_{i-1}^n) 
            - \frac{a\Delta t}{2\Delta x}\bigl(W_{i+1/2}^n - W_{i-1/2}^n\bigr).
\]

Введём параметр $\sigma$ и модифицируем:
\[
u_i^{n+1} = \Bigl[\,\sigma u_i^n 
            + (1-\sigma)\frac{u_{i+1}^n + u_{i-1}^n}{2}\Bigr]
            - \frac{a\Delta t}{\Delta x}\bigl(W_{i+1/2}^n - W_{i-1/2}^n\bigr).
\]

После раскрытия потоков разностное уравнение принимает форму

\begin{equation*}
u_i^{n+1} = \biggl[\,\sigma u_i^n 
            + (1-\sigma)\frac{u_{i+1}^n + u_{i-1}^n}{2}\biggr]
            - a\Delta t\,\frac{u_{i+1}^n - u_{i-1}^n}{2\Delta x}
            + \frac{a^2\Delta t^2}{2}\,
              \frac{u_{i+1}^n - 2u_i^n + u_{i-1}^n}{\Delta x^2}.
\label{eq:modified-expanded}
\end{equation*}
\end{frame}


%=====================================================
\begin{frame}{Анализ ошибки аппроксимации}

Рассмотрев дифференциальное приближение схемы мы можем найти ошибку

\begin{dmath*}
  u_i^{n+1} - u(x_i, t^{n+1}) \approx 
  \input{attachments/scheme-taylor-error.tex} + \ldots
\end{dmath*}

Для повышения порядка разложим производную $u_{xxx}$ компенсируя старшие члены ошибки предыдущих производных.
\end{frame}


%=====================================================
\begin{frame}{Определение весов коррекции}
Аппроксимируем $u_{xxx}$ в узле $i$ на пятиточечном шаблоне:
\[
\left(\frac{\partial^3 u}{\partial x^3}\right)_i
\approx A\,u_{i-2} + B\,u_{i-1} + C\,u_i + D\,u_{i+1} + E\,u_{i+2}.
\]
Подстановка тейлоровских разложений даёт линейную систему относительно $A,\dots,E$:

\vspace{0.3cm}
\[
\input{attachments/taylor-coefficients-system.tex}.
\]
\end{frame}


%=====================================================
\begin{frame}{Итоговая схема повышенного порядка}
Подстановка найденных коэффициентов приводит к

% \small
\begin{dmath*}
u_i^{n+1} =
\input{attachments/scheme-cfl-form-stencil-form.tex},
\end{dmath*}
\normalsize
где $\nu = a\Delta t/\Delta x$ — число Куранта.

Главный член ошибки после коррекции имеет вид
\[
\input{attachments/principal-error-term.tex}.
\]

Итого, порядок интегральной ошибки — третий (при $\Delta t \sim \Delta x$).
\end{frame}


%=====================================================
\begin{frame}{Анализ устойчивости по фон Нейману}
Подстановка гармоники $u_j^n = \lambda^n e^{ij\theta}$ ($\theta = k\Delta x$) даёт
множитель перехода $G_\sigma(\theta,\nu)$. Квадрат модуля $|G|^2$ — полином
от $x = \cos\theta$ и $\nu^2$:
\[
P(x,\nu) = |G_\sigma(\theta,\nu)|^2,\qquad x\in[-1,1].
\]
Схема устойчива тогда и только тогда, когда $\displaystyle\max_{x\in[-1,1]} P(x,\nu) \le 1$.

Граница устойчивости ищется из условий:
\begin{itemize}
\item $P(-1,\nu) = 1$ и $P(1,\nu) = 1$ (граничные волновые числа);
\item $P(x,\nu) = 1$ и $\frac{\partial P(x,\nu)}{\partial x} = 0$ (точки касания). То есть $\operatorname{Disc}_x\bigl(P(x,\nu)-1\bigr)=0$.
\end{itemize}
\end{frame}


%=====================================================
\begin{frame}{Определяющие полиномы и дискриминант}
Граничные условия дают полиномы (по $\nu^2$):
\[
P(1,\nu)-1 \equiv
\input{attachments/left-boundary-stability-condition-polynomial.tex},
\]
\[
P(-1,\nu)-1 =
\input{attachments/right-boundary-stability-condition-polynomial.tex}.
\]

Дискриминант по $x$:
\begin{dmath*}
\operatorname{Disc}_x\bigl(P(x,\nu)-1\bigr) =
\input{attachments/discriminant.tex}
\end{dmath*}

Из этих выражений аналитически находится $\nu_{\max}(\sigma)$.
\end{frame}


%=====================================================
\begin{frame}{Зависимость $\nu_{\max}$ от параметра $\sigma$}
\begin{columns}[T]
\column{0.5\textwidth}
% \left
Поверхность дискриминанта:

% \vspace{-1cm}

\hspace*{-2.5cm}
% \begin{figure}
%     \scalebox{0.45}{\input{attachments/discriminant-surface.pgf}}
% \end{figure}
\scalebox{0.45}{\input{attachments/discriminant-surface.pgf}}
% \input{attachments/discriminant-surface.pgf}
\column{0.5\textwidth}
% \centering
Максимальное число Куранта:
\vspace{0cm}
\vspace*{1cm}
\hspace*{0.5cm}
\scalebox{0.5}{\input{attachments/cfl-bound-vs-parameter.tikz}}
\end{columns}

% $\nu_{\max}(\sigma)$ — минимальное положительное решение $P=1$;
% для $\sigma \in [0.5,\, 1.0]$ значение $\nu_{\max}$ близко к единице.
\end{frame}


%=====================================================
\begin{frame}{Диссипативные и дисперсионные свойства}
\begin{center}
\scalebox{0.75}{\input{attachments/dissipation-dispersion-profile.tikz}}
\end{center}

\vspace{-0.3cm}
\begin{itemize}
\item $|G|\to 1$ при $\theta\to0$ (аппроксимация);
\item Сильная диссипация коротких волн при средних $\nu$;
% \item Относительная фазовая ошибка меняет знак в зависимости от $\nu$.
\end{itemize}
\end{frame}


%=====================================================
\begin{frame}{Порядок сходимости}
\begin{center}
\scalebox{0.8}{\input{attachments/convergence-ge.tikz}}
\end{center}
% \vspace{0.2cm}
% \scalebox{0.8}{\input{attachments/convergence-le-table.tex}}
% \centering
% Интегральная/ ошибка ($\sim \Delta x^3$)

% \vspace{0.2cm}
% \scalebox{0.55}{\input{attachments/convergence-ge.tikz}}

% \vspace{0.2cm}
% \scalebox{0.8}{\input{attachments/convergence-ge-table.tex}}
% \end{minipage}

Теоретические оценки полностью подтверждаются.
\end{frame}


%=====================================================
\begin{frame}{Монотонизация}
Схемы высокого порядка генерируют нефизичные осцилляции (явление Гиббса):
\begin{center}
\scalebox{0.75}{\input{attachments/scheme-comparison-discontinuous.tikz}}
\end{center}
\end{frame}


%=====================================================
\begin{frame}{Метод адаптивной искусственной вязкости}
    \vspace{.2cm}
    Вводится локальное диссипативное слагаемое:
    % Changed for layout adjustment. Same as
    % \(\frac{\partial u}{\partial t} + a\frac{\partial u}{\partial x}
    %   = \frac{\partial}{\partial x}\!\left(\mu(x,t)\frac{\partial u}{\partial x}\right)\).
    \(u'_t + a u'_x = \left(\mu(x,t) u'_x\right)'_x\).
    
    % Алгоритм:
    \begin{multicols}{2}
    \begin{enumerate}
        % \column{0.5\textwidth}
        \vspace{-0.4cm}
        \item Предиктор по базовой схеме
        % высокого порядка
        % \vspace{-0.4cm}
        \item Обнаружение немонотонности
        \vspace{-0.4cm}
        \item Введение вязкости $\mu$ от $\mu_{\min}$ до $\mu_{\max}$
        % (при необходимости $\mu_{\max}$)
        % \vspace{-0.4cm}
        \item Корректор с подобранным полем $\mu$
        % \vspace{-0.4cm}
    \end{enumerate}
    \end{multicols}


    \begin{center}
    \vspace{0.15cm}
    \scalebox{0.7}{
      \input{attachments/adaptive-artificial-viscosity.tikz}
    }
    % \includegraphics[width=\textwidth]{adaptive-artificial-viscosity.tikz}
    \end{center}
% Note: Осцилляции подавлены, но разрыв несколько размывается.
\end{frame}


%=====================================================
\begin{frame}{Монотонизация: гибридная схема}
Комбинируем решения высокого и низкого порядков:
\[
u^{n+1}_i = w_i\, u^{\text{high}}_i + (1-w_i)\, u^{\text{low}}_i,
\]
где $w_i\in[0,1]$ — вес, получаемый из \textbf{датчика осцилляций}.

% \begin{itemize}
% \item датчик кривизны,
% \item датчик вариации на 5‑точечном шаблоне.
% \end{itemize}

\begin{center}
\scalebox{0.75}{\input{attachments/hybrid.tikz}}
\end{center}
\end{frame}


%=====================================================
\begin{frame}{Сравнение способов монотонизации}
  
% Alignment here is a bit finicky
\vspace{0.1cm}
\begin{columns}
\column{0.5\textwidth}
    \textbf{Адаптивная вязкость}
    \begin{itemize}
    % \vspace{-0.2cm}
    % \item[$+$] простая реализация
    \vspace{-0.2cm}
    \item[$+$] устойчива
    \vspace{-0.2cm}
    \item[$-$] размазывает разрыв
    \end{itemize}
\vspace{0.2cm}
\column{0.5\textwidth}
    \textbf{Гибридная схема}
    \begin{itemize}
    \vspace{-0.2cm}
    \item[$+$] крутой профиль разрыва и высокая точность в гладких зонах
    \vspace{-0.2cm}
    \item[$-$] зависит от задачи
    \end{itemize}
\end{columns}

\begin{center}
\scalebox{0.75}{\input{attachments/monotization-comparison.tikz}}
\end{center}
\end{frame}

% \begin{frame}{Сравнение методов монотонизации}
% \vspace{-0.5cm}
% \begin{center}
% \scalebox{0.75}{\input{attachments/monotization-comparison.tikz}}
% \end{center}

% \vspace{-0.3cm}
% \begin{tabular}{@{}p{0.45\textwidth}@{\hspace{0.1\textwidth}}p{0.45\textwidth}@{}}
% \textbf{Адаптивная вязкость} & \textbf{Гибридная схема} \\
% \hline
% \raggedright
% -- размазывает разрыв & -- разрыв остаётся крутым \\
% -- требуется несколько проходов & -- два вычисления потоков \\
% -- устойчива & -- зависит от задачи \\
% \end{tabular}
% \end{frame}


%=====================================================
\begin{frame}{Основные результаты}
\begin{itemize}
\item Построено однопараметрическое семейство схем на базе схемы Лакса--Фридрихса с варьируемыми диссипативными свойствами.
\item Путём коррекции на пятиточечном шаблоне достигнут повышенный порядок аппроксимации.
\item С помощью СКА получены аналитические выражения множителя перехода и точная граница устойчивости $\nu_{\max}(\sigma)$.
\item Предложены и реализованы два метода подавления осцилляций: адаптивная искусственная вязкость и гибридная схема с датчиками кривизны.
% \medskip
% \item \textbf{Схема сочетает высокую точность на гладких решениях и монотонность вблизи разрывов.}
\end{itemize}
\end{frame}


%=====================================================
% Finishing slide
%=====================================================
\begin{frame}
\centering
\Large
Спасибо за внимание!

\bigskip

Вопросы?
\end{frame}

\end{document}
